%-*-latex-*-
%-*-latex-*-
\newcommand\COURSE{ciss350}
\newcommand\ASSESSMENT{q10705}
\newcommand\ASSESSMENTTYPE{Quiz}
\newcommand\POINTS{\textwhite{xxx/xxx}}

\input{myquizpreamble}
\input{yliow}
\input{\COURSE}
%-*-latex-*-

%-*-latex-*-

%-*-latex-*-
\renewcommand\TITLE{\ASSESSMENTTYPE \ \ASSESSMENT}

\renewcommand\EMAIL{}
\renewcommand\AUTHOR{eagraham1@cougars.ccis.edu}

\textwidth=6in
\begin{document}
\topmatterthree

The following is the node for the binary search tree.
\begin{Verbatim}[frame=single,fontsize=\footnotesize]
class BSTNode
{
public:
    int key_;
    BSTNode * parent_;
    BSTNode * left_, * right_;
};
\end{Verbatim}

%------------------------------------------------------------------------------
\nextq
Complete the following function to insert \verb!key! into a BST.
For simplicity, assume that \verb!key! is not in the BST.
\\
\ANSWER
\begin{answercode}
void bst_insert(BSTNode ** p, int key)
{
  if (*p == NULL)
  {
    *p = new BSTNode();
    (*p)->key_ = key;
    (*p)->left_ = NULL;
    (*p)->right_ = NULL;
    return;
  }
  else
  {
    if ((*p)->left_ == NULL)
    {
      (*p)->left_ = new BSTNode();
      (*p)->left_->key_ = key;
      (*p)->left_->left_ = NULL;
      (*p)->left_->right_ = NULL;
      (*p)->left_->parent_ = *p;
    }
    else
    {
      bst_insert(&((*p)->left_), key);
    }
  }
  
  else
  {
    if ((*p)->left_ == NULL)
    {
      (*p)->right_ = new BSTNode();
      (*p)->right_->key_ = key;
      (*p)->right_->left_ = NULL;
      (*p)->right_->right_ = NULL;
      (*p)->right_->parent_ = *p;
    }
    else
    {
      bst_insert(&((*p)->right_), key);
    }
  }
}
\end{answercode}

For the following questions,
to describe this BST
\begin{console}[fontsize=\footnotesize]
            (10)
          /      \
      ( 4)        (25)
     /           /    \
 ( 1)        (12)      (27)
                           \
                            (29)
\end{console}
in text, write
\begin{console}[fontsize=\footnotesize]
[10, [4, 25]]
[4, [1, None]]
[25, [12, 27]]
[1, [None, None]]
[12, [None, None]]
[27, [None, 29]]
[29, [None, None]]
\end{console}

%------------------------------------------------------------------------------
\nextq
What is the height of an empty tree?
\\
\ANSWER
\begin{answercode}
-1
\end{answercode}

%------------------------------------------------------------------------------
\nextq
What is the height of this tree?
\begin{console}[fontsize=\footnotesize]
            (10)
          /      \
      ( 4)        (25)
     /           /    \
 ( 1)        (12)      (27)
                           \
                            (29)
\end{console}
\\
\ANSWER
\begin{answercode}
3
\end{answercode}

%------------------------------------------------------------------------------
\nextq
Let $T$ be a BST with values $1, 4, 10, 12, 25, 27, 29$ with the
\textit{smallest} possible height.
Describe $T$ in text (see above on how to describe a BST in text)
\\
\ANSWER
\begin{answercode}
[12, [4, 27]]
[4, [1, 10]]
[27, [25, 29]]
[1, [None, None]]
[10, [None, None]]
[25, [None, None]]
[29, [None, None]]
\end{answercode}

%------------------------------------------------------------------------------
\nextq
List the least number of key(s) that when deleted from this BST
\begin{console}[fontsize=\footnotesize]
            (10)
          /      \
      ( 4)        (25)
     /           /    \
 ( 1)        (12)      (27)
                           \
                            (29)
\end{console}
will make the resulting BST unbalanced.
If you need to delete more than one key, list the keys in ascending order
and separate by commas.
If there are many ways to make this BST unbalanced,
choose the key(s) that comes first in the dictionary order.
For instance if you think deleting $12, 25$ and $12, 27$ will make the
BST unbalanced, you should enter $12, 25$.
\\
\ANSWER
\begin{answercode}
Delete 1
\end{answercode}

%------------------------------------------------------------------------------
%-*-latex-*-
\newpage
\input{instructions}
\end{document}

